% TOP_PROBLEM: {"id": 3255, "title": "The Two Color Conjecture", "category_id": 3, "category": {"id": 3, "name": "graph_theory", "display_name": "Graph Theory", "description": "Problems involving graphs, networks, and their properties.", "slug": "graph-theory", "order_index": 3, "created_at": "2026-07-31T15:26:25.671Z"}, "status": "open", "problem_number": "OPG-169", "set_id": 12}
% FIRST_POSTED: 2026-10-02
% ATTEMPT_STATUS: unresolved
% UnsolvedMath ID 3255; source catalogue code OPG-169
% !TeX program = lualatex
% Research attempt by GPT-6 Astra Ultra; no claim of a new complete solution.
\documentclass[11pt,a4paper]{article}
\usepackage[margin=25mm]{geometry}
\usepackage{fontspec}
\setmainfont{lmroman10-regular.otf}[BoldFont=lmroman10-bold.otf,
 ItalicFont=lmroman10-italic.otf,BoldItalicFont=lmroman10-bolditalic.otf]
\usepackage{amsmath,amssymb,mathtools}
\usepackage{unicode-math}
\setmathfont{latinmodern-math.otf}
\AtBeginDocument{\let\setminus\smallsetminus}
\usepackage{xurl,hyperref}
\hypersetup{colorlinks=true,urlcolor=blue}
\setcounter{secnumdepth}{0}
\setlength{\parindent}{0pt}
\setlength{\parskip}{5pt}
\setlength{\emergencystretch}{3em}
\begin{document}
\section*{The Two Color Conjecture}\label{3255}
{\small UnsolvedMath ID 3255 · OPG-169 · Graph Theory · OpenGarden}
\subsection{Definitions and mathematical statement}
An oriented graph $D=(V,A)$ is obtained from a finite simple undirected
graph by assigning one direction to each edge. In particular it has no
loops, no repeated arcs, and no pair of opposite arcs. Assume its
underlying undirected graph is planar. A directed cycle is a cyclic
sequence of distinct vertices $v_1,\ldots,v_m$, with $m\ge3$, such that
$v_i v_{i+1}\in A$ for all indices read modulo $m$.

For $X\subseteq V$, the induced subdigraph $D[X]$ retains all arcs whose
ends lie in $X$. It is acyclic if it has no directed cycle. The
Neumann-Lara two-colour conjecture asks whether there is always a
partition
\[
                  V=X_1\sqcup X_2
                  \quad\text{with }D[X_1],D[X_2]\text{ acyclic}.
\]
Empty parts are allowed. Equivalently, there should be a map
$c:V\to\{1,2\}$ such that every directed cycle contains both colours.
The least number of acyclic induced colour classes is called the
dichromatic number, so the requested bound is at most two.

This is a directed condition. A colour class may contain an undirected
cycle if its orientation is not cyclic, and its vertices need not form
an independent set. Conversely, the prohibition of opposite arcs in an
oriented graph is material: a pair of opposite arcs would itself be a
directed cycle of length two. The assertion covers every orientation of
every finite simple planar graph, without any assumption about the
lengths of its directed cycles.
\subsection{Short English statement}
Can the vertices of every planar oriented graph be coloured with two colours so that no directed cycle is monochromatic?
\subsection{Research attempt}
\paragraph{Scope and process.}
This is a GPT-6 Astra Ultra research attempt dated 2 October 2026, using
primary-source browsing and elementary mathematical checks. The canonical
definition above is retained. Results below are restricted results or
reductions, not a claim of a new solution to the entire catalogue question.
No formal proof assistant or external mathematical referee checked this text.


\paragraph{Literature and directed hypotheses.}
Li and Mohar [R1] prove the stronger restricted statement that every
planar digraph of digirth at least four is two-colourable. Their
hypothesis excludes directed triangles, rather than all undirected
triangles. The elementary result below applies to a smaller class but
has a short proof and extends to a general degeneracy bound. Throughout,
there are no opposite arcs; this is essential to the neighbour count.
No identification is made with other conjectures about orientations
forcing large undirected chromatic number.

\paragraph{A degeneracy bound.}
Suppose the underlying graph is $d$-degenerate, meaning that every
nonempty subgraph has a vertex of degree at most $d$. Repeatedly delete
such a vertex and restore vertices in reverse order. Set
$q=\lfloor d/2\rfloor+1$. Inductively colour the already restored
vertices with $q$ colours so that every colour class is acyclic.
When restoring $v$, at most $d$ neighbours are already coloured.
A colour can become invalid at $v$ only if $v$ has both an incoming and
an outgoing neighbour of that colour: any newly created directed cycle
must pass through $v$ and use both types of arc. Those two neighbours
are distinct because opposite arcs are prohibited. If all $q$ colours
were invalid, there would therefore be at least $2q>d$ coloured
neighbours, a contradiction. Choose a colour missing one of the two
types. This proves
\[
             \vec\chi(D)\le \lfloor d/2\rfloor+1.
\]
The condition is sufficient, without requiring a search for directed
paths between the neighbours.

Euler's inequality gives $e\le3v-6$ for a simple planar graph with at
least three vertices; hence every such graph is 5-degenerate, including
its small subgraphs. The construction yields three acyclic colours for
all planar orientations. For triangle-free planar graphs,
$e\le2v-4$ gives 3-degeneracy and hence two colours. Forest components
and subgraphs with fewer than three vertices satisfy the same
degeneracy conclusion directly. This settles a genuine infinite
subclass with arbitrary orientations.

\paragraph{Necessary structure of a counterexample.}
Suppose $D$ were vertex-minimal among planar orientations not admitting
two acyclic colours. Every vertex must have at least two incoming and
two outgoing neighbours. Indeed, delete $v$ and colour the remainder.
If $d^-(v)\le1$, one colour appears on no incoming neighbour, so assigning
that colour to $v$ creates no monochromatic directed cycle. The same
argument applies to $d^+(v)\le1$. Thus, with $n$ vertices and $m$ arcs,
\[
                  \delta(D_{\rm und})\ge4,\qquad m\ge2n.
\]
Moreover $D$ is strongly connected: otherwise its strongly connected
components are proper induced subdigraphs and can each be two-coloured;
every directed cycle lies within one component, so their colourings
combine. Its underlying graph consequently has no bridge. A bridge
with only one orientation cannot allow directed paths in both
directions between its two sides.

Embed the connected underlying graph in the plane. Every face boundary
walk has length at least three, since the graph is simple and
bridgeless. With $f=m-n+2$ faces, counting boundary lengths gives
\[
  \sum_F(4-|\partial F|)=4f-2m=8-4n+2m\ge8.
\]
A triangular face contributes one, a quadrilateral zero, and a larger
face a negative quantity. Therefore every such minimal counterexample
has at least eight triangular faces. This counts undirected facial
triangles; it does not assert that all eight have cyclic orientations.
It is a restriction on possible counterexamples, not a contradiction.

\paragraph{Verification and the stalled extension.}
The certificate [R2] checks every orientation of $K_4$ by testing all
vertex bipartitions for acyclicity. This is a small sanity check on the
directed condition, not an exhaustive test of planar graphs. The proof
above treats arbitrarily large triangle-free planar graphs.
At a vertex of underlying degree four or five, two colours can each
occur on both an incoming and an outgoing neighbour. The local
extension rule then stops: deciding whether those neighbours are joined
by monochromatic directed paths depends on the rest of the graph.
The face count does not remove that obstruction.

\paragraph{Final status and first remaining gap.}
\textbf{UNRESOLVED.} We obtain the general three-colour bound, prove
the triangle-free underlying case with two colours, and constrain a
minimal counterexample. An extension argument handling the interaction
of directed triangles and those monochromatic paths is still missing
for arbitrary planar orientations.

\subsection{Sources}
\begin{itemize}
\item[S1] ULAM.AI and the UnsolvedMath Contributors, supplied problems.json, record 3255 (OPG-169), OpenGarden collection. Catalogue identity and source transcription; CC BY 4.0.
\url{https://huggingface.co/datasets/ulamai/UnsolvedMath}.
\item[S2] Open Problem Garden, The Two Color Conjecture. Original problem and discussion; the mathematical formulation below is expanded from the supplied source record.
\url{http://www.openproblemgarden.org/op/partitioning_planar_digraphs}.

\item[R1] Zhentao Li and Bojan Mohar, \emph{Planar digraphs of digirth
four are 2-colourable}, 2016 preprint. Primary abstract checked
2 October 2026. \url{https://arxiv.org/abs/1606.06114}.
\item[R2] Reproducible finite checks:
\path{scripts/check_attempt_batch03_colouring_2026_10_02.py}, section 3255.

\end{itemize}
\end{document}
